% 2009 年辽宁卷文科第 20 题：甲、乙两分厂零件内径尺寸的频数分布表。
% 来源 PDF：普通高考/2009/2009辽宁文.pdf，SHA256 99623efa457632f26bc5d5775293bf9fd13a56add984179360e30698afc7cb73
% 原页：第 2 页（A3 三栏）第 2 栏上部，印刷页 540。
% 重绘依据：两表各 7 个分组，组距 0.04，分组与频数逐字照录原页；第三个分组原页印作
%   [29.84,29.98)（两表相同），与前后分组重叠，照录并登记问题，未擅改。
\begin{tikzpicture}[semithick, every node/.style={inner sep=2pt, font=\small}]
  % 甲厂
  \node at (1.45,0.35) {甲厂};
  \draw (0,0) rectangle (4.0,-4.55);
  \draw (2.9,0) -- (2.9,-4.55);
  \draw (0,-0.7) -- (4.0,-0.7);
  \foreach \y in {-1.25,-1.8,-2.35,-2.9,-3.45,-4.0} { \draw (0,\y) -- (4.0,\y); }
  \node at (1.45,-0.35) {分组};
  \node at (3.45,-0.35) {频数};
  \node at (1.45,-0.975) {$[29.86,\,29.90)$};
  \node at (1.45,-1.525) {$[29.90,\,29.94)$};
  \node at (1.45,-2.075) {$[29.84,\,29.98)$};
  \node at (1.45,-2.625) {$[29.98,\,30.02)$};
  \node at (1.45,-3.175) {$[30.02,\,30.06)$};
  \node at (1.45,-3.725) {$[30.06,\,30.10)$};
  \node at (1.45,-4.275) {$[30.10,\,30.14)$};
  \node at (3.45,-0.975) {$12$};
  \node at (3.45,-1.525) {$63$};
  \node at (3.45,-2.075) {$86$};
  \node at (3.45,-2.625) {$182$};
  \node at (3.45,-3.175) {$92$};
  \node at (3.45,-3.725) {$61$};
  \node at (3.45,-4.275) {$4$};
  % 乙厂
  \begin{scope}[shift={(5.2,0)}]
    \node at (1.45,0.35) {乙厂};
    \draw (0,0) rectangle (4.0,-4.55);
    \draw (2.9,0) -- (2.9,-4.55);
    \draw (0,-0.7) -- (4.0,-0.7);
    \foreach \y in {-1.25,-1.8,-2.35,-2.9,-3.45,-4.0} { \draw (0,\y) -- (4.0,\y); }
    \node at (1.45,-0.35) {分组};
    \node at (3.45,-0.35) {频数};
    \node at (1.45,-0.975) {$[29.86,\,29.90)$};
    \node at (1.45,-1.525) {$[29.90,\,29.94)$};
    \node at (1.45,-2.075) {$[29.84,\,29.98)$};
    \node at (1.45,-2.625) {$[29.98,\,30.02)$};
    \node at (1.45,-3.175) {$[30.02,\,30.06)$};
    \node at (1.45,-3.725) {$[30.06,\,30.10)$};
    \node at (1.45,-4.275) {$[30.10,\,30.14)$};
    \node at (3.45,-0.975) {$29$};
    \node at (3.45,-1.525) {$71$};
    \node at (3.45,-2.075) {$85$};
    \node at (3.45,-2.625) {$159$};
    \node at (3.45,-3.175) {$76$};
    \node at (3.45,-3.725) {$62$};
    \node at (3.45,-4.275) {$18$};
  \end{scope}
\end{tikzpicture}
